\documentclass[10pt,a4paper]{article}
\usepackage[top=1.4cm,bottom=1.4cm,left=2.0cm,right=2.0cm]{geometry}
\usepackage{booktabs}
\usepackage{array}
\usepackage[dvipsnames]{xcolor}
\usepackage{parskip}
\pagestyle{empty}

\newcommand{\good}[1]{\textcolor{OliveGreen}{#1}}
\newcommand{\bad}[1]{\textcolor{BrickRed}{#1}}

\begin{document}

\begin{center}
{\Large\bfseries Correction Memo --- VANI Speech-Recognition Results}\\[2pt]
{\small 12 July 2026, addendum 13 July \quad$\cdot$\quad replaces the fine-tuning figures reported before 10 July 2026}
\end{center}

\textbf{In one paragraph.} While building a new evaluation tool, I discovered that the numbers we had
reported for our fine-tuned speech-recognition models were distorted by three hidden mistakes in how the
models were tested and trained --- not in the models themselves. I have fixed all three, re-measured
everything, and re-checked every figure with automated scripts. The corrected results change our
conclusion: custom training (``fine-tuning'') is worth it for only two of our seven languages; for the
other five, an off-the-shelf model called SeamlessM4T is simply better, and the system now uses it.
The corrected story is more defensible than the original one.

\textbf{How we measure.} Accuracy is reported as \emph{Word Error Rate} (WER): the percentage of words
the system gets wrong --- \textbf{lower is better}. A WER of 20\% means roughly one word in five is wrong.

\textbf{Mistake 1 --- we compared against the wrong ``before''.} Every fine-tuned model was compared
against an un-tuned baseline to show improvement. Due to a mislabelled file, the baseline we actually
loaded was a \emph{smaller, weaker} model than the one named in the report. A weaker ``before'' makes
any ``after'' look more impressive, so every improvement we reported was overstated --- for Urdu, the
real gain almost vanishes (from $-4.6$ points to $-1.4$).

\textbf{Mistake 2 --- Mandarin was scored by a rule that cannot work for Chinese.} WER counts words by
the spaces between them, but Chinese is written without spaces. Our reference texts happened to contain
spaces between every character while the model's output (correctly) had none --- so a near-perfect
transcript was scored as $\sim$100\% wrong. The headline result ``Mandarin improved from 100\% error to
16\%'' was entirely this scoring quirk. Scored properly (character by character), the un-tuned model was
already at 10.99\% and our fine-tuning actually made Mandarin slightly \emph{worse} (14.22\%).

\textbf{Corrected results} (100 test recordings per language; best score in bold; the last column is what
the system now actually uses):

\begin{center}
\small
\begin{tabular}{@{}l r r r r l@{}}
\toprule
 & \multicolumn{1}{c}{Old report} & \multicolumn{3}{c}{Corrected WER (\%)} & \\
\cmidrule(lr){2-2}\cmidrule(lr){3-5}
Language & before $\rightarrow$ after & un-tuned & our fine-tune & SeamlessM4T & System now uses \\
\midrule
Punjabi  & $105.8 \rightarrow 55.7$ & 77.6 & 57.4 & \textbf{19.8} & SeamlessM4T \\
Nepali   & $94.6 \rightarrow 49.2$  & 88.9 & 50.9 & \textbf{28.5} & SeamlessM4T \\
Hindi    & $30.3 \rightarrow 19.8$  & 26.3 & 19.8 & \textbf{15.4} & SeamlessM4T \\
Urdu     & $24.4 \rightarrow 19.8$  & 21.2 & 19.8 & \textbf{16.9} & SeamlessM4T \\
Mandarin & $100.0 \rightarrow 16.0$ & \textbf{11.0} & \bad{14.2} & 11.7 & SeamlessM4T \\
Pashto   & $94.2 \rightarrow 38.6$  & 89.8 & \textbf{38.6} & 44.4 & \good{our fine-tune} \\
Kashmiri & $98.6 \rightarrow 103.6$ & 96.9 & \textbf{74.0} & unsupported & \good{our fine-tune} \\
\bottomrule
\end{tabular}
\end{center}

\textbf{What the corrected numbers say.} Fine-tuning helps most where the base model knows the language
least (Pashto, Kashmiri) and can even hurt where the base model is already strong (Mandarin). We also
verified, on noisy radio-style audio, that SeamlessM4T's advantage \emph{grows} as conditions worsen ---
so the choice holds in the field, not just in the lab. This ``pick the right engine per language''
result is a stronger contribution than the original ``fine-tuning always wins'' claim.

\textbf{Addendum, 13 July --- Mistake 3, found in the training recipe.} When we fine-tuned SeamlessM4T,
every practice example was accidentally tagged with the wrong language label (``English'') due to a
subtle software default. The model still learned to \emph{transcribe} correctly --- which is why nothing
looked wrong --- but the mislabelling quietly broke its ability to \emph{translate}, and we had
published that breakage as a genuine scientific finding. Re-training one language (Pashto) with the
correct labels restored translation almost completely (quality score 2.1 $\rightarrow$ 37.6, near the
un-tuned 40.2) while transcription stayed identical --- proving the label was the culprit. That earlier
``finding'' is withdrawn. No deployment decision changes.

\textbf{Why you can trust the corrected numbers.} All three mistakes were found in-house, by our own
checks. The old results are archived, every re-run keeps its raw transcripts so scores can be re-checked
without re-running models, and two automated scripts now re-verify all 380 figures in the reports
against the raw data after any edit. The corrected system configuration is deployed and tested end-to-end.

\vspace{2pt}
--- VANI project, 13 July 2026

\end{document}
